\chapter{Formulação do Problema e Métricas}\label{cap:proposta}

\section{Definição do Problema}
O problema central abordado neste trabalho é a rigidez das políticas de compactação tradicionais em LSM-trees diante da variabilidade da Memória Persistente. Em motores de armazenamento convencionais, os gatilhos de compactação são puramente reativos e baseados em limiares estáticos (por exemplo, número de arquivos no nível L0 ou tamanho total de um nível). Essa abordagem falha em sistemas baseados em PMem por dois motivos principais:

\begin{itemize}
    \item Sensibilidade à Latência de Escrita: A PMem, embora rápida, possui uma largura de banda de escrita assimétrica e limitada se comparada à leitura. Compactações agressivas durante picos de carga de escrita geram contenção no controlador de memória, levando ao fenômeno de write stall (parada total da ingestão de dados). 
    \item Ignorância do Estado do Sistema: As regras estáticas não consideram o custo marginal de uma compactação em um dado momento. Uma compactação iniciada durante uma saturação de banda pode ser mais prejudicial ao desempenho global do que o aumento temporário na amplificação de leitura causado pelo adiamento dessa mesma compactação.
\end{itemize}

Dessa forma, o problema reside em como transformar o processo de compactação de um "mal necessário" executado às cegas em um processo consciente do hardware e da carga de trabalho, que maximize a vazão (throughput) sem comprometer a integridade da estrutura da árvore LSM.

\section{Proposta de Solução}

Propõe-se o Adaptive Compaction Scheduler (ACS), um escalonador de compactação orientado a feedback. O ACS substitui os gatilhos fixos por uma função de utilidade dinâmica. Ele monitora métricas de hardware (latência de persist-barriers) e métricas da árvore LSM para decidir o grau de concorrência e a prioridade das tarefas, adiando compactações não críticas durante períodos de pressão extrema na PMem e acelerando-as quando o sistema detecta ociosidade.
